Este trabalho teve como objetivo propor uma solução computacional para o problema de alocação semestral de grades horárias em Instituições de Ensino Superior, por meio do desenvolvimento do IGRIS, um sistema modular de suporte à decisão baseado em Algoritmos Genéticos. A proposta partiu da necessidade de automatizar um processo tradicionalmente conduzido de forma manual, caracterizado por elevada complexidade combinatória e pela necessidade de considerar simultaneamente disciplinas, turmas, docentes, disponibilidades e restrições acadêmicas.

A solução desenvolvida demonstrou que o uso de Algoritmos Genéticos é uma alternativa tecnicamente viável para apoiar a geração automatizada de grades horárias acadêmicas. Por meio da representação das alocações como soluções candidatas, foi possível aplicar operadores evolutivos, como seleção, \textit{crossover}, mutação e elitismo, com o objetivo de explorar diferentes combinações de horários e buscar arranjos mais adequados às restrições do problema. Dessa forma, o IGRIS contribui para reduzir o esforço manual envolvido na montagem das grades e para apoiar os coordenadores acadêmicos no processo de tomada de decisão.

A arquitetura modular adotada também se mostrou adequada aos objetivos do trabalho, pois permitiu separar as responsabilidades entre interface, persistência, serviços de apoio e módulos de otimização. Essa organização favorece a manutenção do sistema, a evolução independente dos componentes e a inclusão de diferentes implementações algorítmicas dentro de uma mesma estrutura de integração. Com isso, o IGRIS não se limita a uma única estratégia de geração de grades, mas oferece uma base para avaliar e comparar diferentes abordagens aplicadas ao \textit{University Course Timetabling Problem}.

No desenvolvimento do sistema, a modelagem das entidades acadêmicas e dos fluxos de operação foi essencial para aproximar o problema real de sua representação computacional. A definição de turmas, disciplinas, docentes, disponibilidades, matrizes curriculares e vínculos permitiu estruturar os dados necessários ao funcionamento dos algoritmos. Essa etapa foi fundamental, pois a qualidade da solução gerada depende diretamente da forma como as informações acadêmicas são organizadas e disponibilizadas para processamento.

As implementações dos algoritmos Morpheus, Sunny e Rocket permitiram avaliar diferentes decisões de projeto aplicadas ao mesmo domínio. Embora os três utilizem princípios fundamentais dos Algoritmos Genéticos, eles diferem em aspectos como organização interna dos dados, tratamento de conflitos, dinâmica de execução, uso de parada antecipada, avaliação da população e consumo de recursos computacionais. Nos cenários analisados, o Rocket apresentou o melhor desempenho operacional, com menores valores de latência média, percentil 95 e consumo de memória RAM. O Morpheus manteve desempenho superior ao Sunny nas métricas operacionais agregadas, enquanto o Sunny apresentou taxa de sucesso elevada e comportamento competitivo em determinados recortes de parametrização, ainda que com maior uso de memória e maior número médio de gerações no conjunto geral dos testes.

Os resultados obtidos indicam que a comparação entre as implementações deve ser compreendida principalmente sob a perspectiva de eficiência computacional. Embora Rocket e Morpheus tenham apresentado melhor desempenho em tempo de resposta, consumo de recursos e score final em relação ao Sunny, os três algoritmos mantiveram erro médio final equivalente nos cenários avaliados. Assim, os dados não permitem afirmar que uma das implementações produza grades qualitativamente superiores em todos os aspectos, mas sustentam que o Rocket alcançou resultados equivalentes com menor custo operacional no protocolo experimental adotado, seguido pelo Morpheus em relação ao Sunny.

Também foi possível observar que a automação proposta não elimina a necessidade de validação humana. O IGRIS atua como uma ferramenta de apoio à decisão, oferecendo soluções candidatas, reduzindo o esforço operacional e organizando o processo de geração de grades. Entretanto, a análise final ainda cabe aos responsáveis acadêmicos, especialmente porque algumas restrições institucionais podem não estar completamente formalizadas no sistema ou depender de critérios administrativos específicos.

A interface web desenvolvida contribui para tornar a solução utilizável no contexto institucional, permitindo o cadastro de dados, a configuração dos parâmetros e a visualização dos resultados de forma organizada. Esse aspecto reforça o caráter aplicado do trabalho, pois a efetividade de uma solução computacional para gestão acadêmica não depende apenas do algoritmo de otimização, mas também da clareza da interação com o usuário, da consistência dos dados e da integração com o fluxo de trabalho dos coordenadores.

Apesar dos resultados alcançados, algumas limitações devem ser consideradas. Os testes realizados não contemplam todos os possíveis cenários acadêmicos, e os relatórios de benchmark não detalham integralmente aspectos como ambiente de execução, controle de aleatoriedade, decomposição dos erros finais e aplicação de critérios estatísticos inferenciais. Além disso, os relatórios do Rocket foram gerados separadamente dos relatórios de Morpheus e Sunny, ainda que com o mesmo perfil de \textit{payload} e os mesmos cenários declarados, o que recomenda cautela adicional na interpretação comparativa. Essas limitações indicam a necessidade de avaliações complementares com maior diversidade de dados, maior controle experimental e registro mais detalhado das condições de execução.

Como trabalhos futuros, recomenda-se ampliar os experimentos com bases acadêmicas de diferentes tamanhos e níveis de complexidade, além de decompor o erro médio por tipo de inconsistência, como conflito de professor, conflito de turma, demanda não alocada e incompatibilidade de disponibilidade. Também se recomenda incorporar mecanismos de observabilidade, como logs estruturados, métricas detalhadas de execução, profiling de memória, rastreamento de alocação, logs de Garbage Collector quando aplicável e registro das sementes aleatórias utilizadas, a fim de aumentar a rastreabilidade e a reprodutibilidade dos resultados.

Outra possibilidade de continuidade consiste na inclusão de novos módulos de otimização, baseados em técnicas como Simulated Annealing, Busca Tabu, Otimização por Enxame de Partículas ou abordagens híbridas. A arquitetura modular do IGRIS favorece essa expansão, permitindo que diferentes estratégias sejam integradas, executadas e comparadas conforme o perfil do problema acadêmico analisado.

Conclui-se, portanto, que o objetivo do trabalho foi alcançado, uma vez que foi proposta e desenvolvida uma solução computacional modular para apoiar a geração automatizada de grades horárias em Instituições de Ensino Superior. O IGRIS demonstrou capacidade de representar o problema de alocação acadêmica, aplicar Algoritmos Genéticos à geração de soluções candidatas, reduzir o esforço manual dos coordenadores e permitir a avaliação comparativa entre diferentes implementações evolutivas. Embora existam limitações experimentais e oportunidades de aprimoramento, os resultados obtidos indicam que a solução proposta constitui uma base consistente para a evolução de ferramentas computacionais de apoio à gestão acadêmica.
